\begin{table}[H]
\centering
\renewcommand{\arraystretch}{1.25}
\captionsetup{justification=raggedright,singlelinecheck=false}

\begin{tabular}{|p{1cm}|p{6cm}|p{6cm}|}
\hline
\textbf{Skala} &
\textbf{Definisi Kemungkinan (\textit{Likelihood})} &
\textbf{Definisi Dampak (\textit{Impact}) pada Penelitian} \\ \hline

5 & Hampir Pasti (>90\% terjadi berulang kali). &
Kritis: Sistem gagal total dan tujuan penelitian tidak tercapai. \\ \hline

4 & Sering (>60\% kemungkinan terjadi). &
Berat: Fitur inti tidak berfungsi (misalnya YAML selalu error) dan membutuhkan perbaikan besar pada model. \\ \hline

3 & Mungkin (50\% terjadi sewaktu-waktu). &
Sedang: Penurunan performa (misalnya respon lambat), namun fungsi utama tetap berjalan. \\ \hline

2 & Jarang (<30\% terjadi). &
Ringan: Gangguan minor yang dapat diperbaiki manual dan tidak mempengaruhi hasil utama penelitian. \\ \hline

1 & Sangat Jarang (<10\% terjadi). &
Tidak Signifikan: Tidak mempengaruhi kesimpulan atau hasil eksperimen. \\ \hline

\end{tabular}
\caption{Dimensi Skala Penilaian Risiko: \textit{Likelihood} dan \textit{Impact}}
\label{tbl:risk_scale}
\end{table}